% F1: About my research. 1 to 1.5 pages of prose, no bullet lists. Cover:
%   - MS or PhD student; your advisor; the research you have done so far
%   - the topics and questions you want to work on, and why they suit your
%     skills and interests
%   - what you know about the state of your field and its open problems, and
%     who would benefit if they were solved
%   - what you want to get out of this course
% Cite at least 3 papers you already know, e.g. \cite{key1,key2,key3}, using
% the keys from references.bib.
% AI use: you may draft this from your own bullet notes. If you do, put the
% notes in appendices/my_material.tex.

\section{About My Research}

The small satellite space has been distrupted significantly by deorbit and orbital debris requirements. Most missions are now required to deorbit in the range of 5 years. In addition, orbital crowding in low-Earth orbit (LEO) is driving spaccraft to higher or lower orbits. In either case, the current environment is broadening the number of missions in which some method of propulsion is needed to artifically limit lifespan.  

Propulsion systems on small satellites are difficult to build: there are many factors which fundamentally limit the performance of different systems. Electric propulsion systems are power-limited by the geometry of the satellite, and thus are religated to small, low total impulse systems. The flight history of these systems is similarly lacking. Bipropellant or monopropellant systems have significant challenges in terms of managing heat, propellant mass, health and safety requirements, and overall system complexity. Some work has even considered solar/drag sails to meet orbit raising or lowering requirements. Resistojets and warm-gas type thrusters are uniquely suited to small satellites as they run in power ranges which are acheivable (20-50W), can use a variety of safe propellants, and can be constructed with minimal complexity or control challenges.

A satellite is not entirely constituted by its propulsion system, however. Propulsion development must be handled in tandem with other subsystems. Many questions arrive once one begins thinking through the process of propulsion integration. For a non-vectored thrust, how can the satellite most efficiently point itself to handle a burn? How will subsystems on the satellite bus handle a large, potentially long-term power draw of a propulsion system? How can energy be suppilied to the propulsion system consistently while accounting for eclipse conditions? What are the most effective sensors and measures to determine the system state of the propulsion system?

This project focuses on the development of a water-based propulsion system which uses a novel phase-change based method to maximize total energy collected over an orbit. Considering the lack of conduction into space, the "thermal battery" (a large heatsink made of phase-change material) included with the system allows the satellite to slowly move energy into the propulsion system when available (i.e., in sunlight). I am under obligations from the original inventor of the propulsion system to be brief with some description of the thruster in regards to some of the internal details, construction techniques, and materials used. The goal of this work is to characterize this propulsion system and answer the above questions for a hypothetical mission.

I am a MS student in signal processing, and my advisor is Dr. James McNames. I also have Andrew Greenberg assisting me with the technical aspects. I performed a research review of small-satellite propulsion technologies in my undergraduate degree as part of a literature review program with the Oregon Space Grant Consortium. Part of this work is in conjunction with the Amateur Satellite Corporation (AMSAT) who is handling some of the test stand development. I am working with two other engineers who are currently building a vacuum test chamber where we'll collect data in Tuscon, Arizona. 

Over the summer of 2026, I've spent time working on modeling some of the thermal aspects of the propulsion system. Specifically, I was attempting to use MATLAB to gain an understanding of heat flow out of the PCM and into the water so that we'd have a model to validate against when taking measurements. In addition, I've created a printed circuit board (PCB) which contains a processor, heater and solenoid drive ciruits, and a variety of system sensors (temperature, pressure, etc.). I have started Fall term working on firmware tasks for this board like data transmission, error handling, and basic control over the heater and thruster. 
